\documentclass[border=2pt]{standalone}
\usepackage{itikz}


\begin{document}
\pgfplotsset{%
    axissty/.style={
        % 光滑曲线
        % smooth,
        % 
        % 用于两个坐标系叠加
        % set layers,
        % 
        % 宽度
        % widht = 10cm,
        % 高度
        % height = 10cm,
        % 收缩比例
        % scale = 1,
        % 
        % x , y坐标轴比例一致
        scale only axis,
        % 
        % 散步图时,自动收缩坐标轴
        enlargelimits=0.2, 
        % 
        % 风格辅助线
        % grid = major,
        % grid style = {color=gray,opacity=0.1,line width =0.1pt},
        % 
        % 图标题
        % title = {$\lambda=1.5,c=0.3,c_r=0.07 $},
        % title style = { align = left ,font = \footnotesize},
        % 
        % 隐藏坐标系；
        % hide axis
        % 坐标轴设置 box | top | middle | center | bottom | none | left | right
        % axis x line = bottom,
        % axis y line = left,
        % 
        % 设置坐标轴方向
        % x dir=reverse,
        % y dir=reverse,
        % 
        % 设置轴偏移
        % axis x line shift = -3,
        % axis y line shift = -3,
        % axis line shift = 6.5pt,
        % 
        % 坐标轴样式
        % x axis line style = {-Latex[],blue,shorten > = -10pt},
        % y axis line style = {-Latex[],red,shorten > = -20pt},
        % axis line style = {-Latex[]},
        % 
        % 设置轴刻度的跳跃值 crunch | parallel | none 
        % axis x discontinuity = crunch ,
        % axis y discontinuity = parallel ,
        % 
        % 坐标轴标签
        % xlabel = x-label,
        % ylabel = y-label,
        % 
        % 坐标轴标签样式 lower | upper | bottom | top | left | right | default
        % ylabel style = { align = left,red,font = \footnotesize},
        % label style = { align = left,red,font = \footnotesize},
        % 
        % 坐标轴上下限
        % xmin = 0, 
        % xmax = 31.1,
        % ymin = 0, 
        % ymax= 7,
        % 
        % 坐标轴刻度
        % xtick = {\xmin,10,12,\xmax},
        % ytick = {\ymin,2.2,1,\ymax},
        % 
        % 坐标轴显示的刻度
        % xtickmin = 4,
        % xtickmax = 5,
        % 
        % 坐标轴子刻度
        % minor xtick={0.5,2.4,3.6,4.8},
        % minor ytick={2.7,3.5,3.7,4.1},
        % 
        % 
        % 设置轴刻度位置 inside | center | outside
        % xtick align = outside,
        % ytick align = outside,
        % tick align = outside,
        % 
        % 刻度距离
        % xtick distance = 1,
        % ytick distance = 1,
        % 
        % 坐标轴刻度标签
        % xticklabels = {0,$a$,1.222,$\gamma_i$},
        % yticklabels = {0,$ \lambda c $,1.08,$v_i$},
        % 
        % 坐标轴标签样式设置
        % x tick label style = {/pgf/number format/fixed zerofill,rotate = 90},
        % y tick label style = {/pgf/number format/precision=2},
        % tick label style = {/pgf/number format/precision=2,rotate = 90,},
        %
        % 
        % 图例的设置
        % legend style = {
        % 	% 设置图例文字对齐方式
        % 	legend cell align = left,
        % 	%设置图例水平显示的例数量
        % 	legend columns = 1,
        %     % 设置图例中 小图的对齐位置
        % 	cells = {anchor= east},
        % 	% 设置图例锚点
        % 	anchor = north east,
        % 	% 图例的位置
        %     % at = {($(current bounding box.north west) - (0mm, 0mm)$) },
        % 	% 设置图例位置 
        % 	legend pos= outer north east,
        %     },
        % 
        % 3D 图 或热力图相关
        % colorbar: horizontal | left | right 
        % colorbar horizontal, % 必须配合算子参数 scatter or mesh 使用
        % blackwhite | greenyellow | violet | PuBu-9  (see pgf_plots.pdf, P420)
        % colormap/PuBu-9,
        % colorbar/width=3mm,
        % colorbar style={
        %     at={($(plot.north) + (0,10mm) $)},
        %     anchor=south,
        %     % right | left | upper | lower
        %     ticklabel pos= right,
        % },
        % colorbar style={y dir=reverse},
        % 
        % 3D x轴60度, y轴30度
        % view={45}{45},
    }
}



\begin{tikzpicture}[even odd rule]
\begin{axis}[axis lines* = left]


\addplot+ [raw gnuplot, no markers,name path= g]%
        gnuplot {%
        set xrange [ 0 : 25 ];
        set yrange [ 0 : 7 ];
        plot sample [*:6] x, [6:8] cos(x), [8:10] (x-8)**2;
    };




% \addplot3+ [ no markers,
%         name path= g,
%         surf,
%         domain = -500:500,
%         y domain = -400:400, 
%         z buffer=sort,]%
%      {sin(sqrt((x-20.)^2+(y-20.)^2))/sqrt((x-20.)^2+(y-20.)^2) };






% \fun[0:24][0:5][F2] { y - 4* (1 - exp(-(6-y)*(23-x)/18)) +3}



% \fill [name intersections={of=F2 and g, name = i,total = \n},fill = red] ;



\end{axis}
\end{tikzpicture}
\end{document}